\section{学习率调度}

\subsection{学习率的重要性}

学习率是训练中最重要的超参数之一。选择不当会导致灾难性后果：

\begin{figure}[H]
\centering
\includegraphics[width=0.95\textwidth]{slides-images/slide-050.jpg}
\caption{不同学习率的效果：过高导致发散，过低导致收敛极慢，适中的学习率能快速稳定收敛\protect\footnotemark}
\end{figure}
\footnotetext{Slides 第54页。}

\begin{itemize}
    \item \textbf{过高}：损失不降反升，模型``跳出''损失景观
    \item \textbf{偏高}：损失无法继续下降，在最优点附近震荡
    \item \textbf{过低}：收敛极其缓慢，浪费计算资源
    \item \textbf{适中}：快速下降后持续改进
\end{itemize}

\subsection{学习率衰减策略}

训练过程中\textbf{动态调整}学习率是现代深度学习的标准做法。几种常见策略：

\textbf{1. 阶梯式衰减}（Step Decay）：每经过固定迭代次数，将学习率乘以一个因子（如 0.1）。训练 ResNet 时常用此策略。

\begin{figure}[H]
\centering
\includegraphics[width=0.95\textwidth]{slides-images/slide-052.jpg}
\caption{阶梯式学习率衰减：在固定时间点将学习率降为 1/10\protect\footnotemark}
\end{figure}
\footnotetext{Slides 第56页。}

\textbf{2. 余弦退火}（Cosine Annealing）：学习率按半个余弦周期从最大值平滑降到零：

$$\alpha_t = \frac{\alpha_0}{2}\left(1 + \cos\left(\frac{t \cdot \pi}{T}\right)\right)$$

\begin{figure}[H]
\centering
\includegraphics[width=0.95\textwidth]{slides-images/slide-053.jpg}
\caption{余弦退火学习率调度：从最大值平滑降至零\protect\footnotemark}
\end{figure}
\footnotetext{Slides 第57页。}

\textbf{3. 线性衰减}和\textbf{逆平方根衰减}等其他方案。

\textbf{4. 线性预热}（Linear Warmup）：训练初期先从很小的学习率线性增长到目标值，然后再按照某种策略衰减。这可以避免初始时不稳定的大步长。

\begin{figure}[H]
\centering
\includegraphics[width=0.95\textwidth]{slides-images/slide-055.jpg}
\caption{线性预热 + 余弦退火：先线性增长到最大学习率，然后按余弦曲线衰减\protect\footnotemark}
\end{figure}
\footnotetext{Slides 第59页。}

\begin{importantbox}{线性缩放规则}
经验法则：如果将 batch size 增大 $n$ 倍，学习率也应相应增大 $n$ 倍。直觉是：更大的 batch 提供更准确的梯度估计，因此可以``迈更大的步''。这一规则在许多实际场景中被证明有效。
\end{importantbox}

\subsection{本章小结}

学习率是最关键的超参数。实践中几乎所有现代模型都使用学习率调度：先预热后衰减。余弦退火 + 线性预热是最流行的组合之一。batch size 增大时应按比例增大学习率。

%% ========== 二阶优化 ==========

